\documentclass[14pt,a4paper,oneside]{extbook}
\usepackage[utf8]{vietnam}
\usepackage{amsmath,amssymb,mathrsfs,fancyhdr,enumerate,multirow,makecell}
\usepackage{tkz-euclide,tikz-3dplot,tikz,tkz-tab}
\usetikzlibrary{arrows.meta,calc,intersections,angles,quotes,backgrounds,shapes.geometric,patterns,through,shadows,positioning}
\usepackage{yhmath}
\usepackage[top=0.7cm, bottom=0.7cm, left=1.0cm, right=1.0cm, headheight=16pt, headsep=3mm, footskip=5mm]{geometry}
\usepackage[hidelinks,unicode]{hyperref}
\usepackage{qrcode}
\usepackage[many]{tcolorbox}
\usepackage{fontawesome}

\pagestyle{fancy}
\fancyhf{}
\fancyhead[L]{\footnotesize\bfseries\color{teal!85!black}\faBook~SỔ TAY KIẾN THỨC TOÁN 4}
\fancyhead[R]{\footnotesize\bfseries\color{blue!85!black}Lớp Toán Thầy Lĩnh -- loptoanthaylinh.vn}
\fancyfoot[C]{\footnotesize\bfseries\color{gray!80!black}-- Sổ tay kiến thức: Bài 58. So sánh hai phân số --}
\renewcommand{\headrulewidth}{0.4pt}

\begin{document}

% --- BANNER TIÊU ĐỀ BÀI HỌC ---
\noindent
\begin{tcolorbox}[
  enhanced,
  colback=teal!8!white,
  colframe=teal!80!black,
  arc=3mm, boxrule=1pt,
  left=3mm, right=3mm, top=2mm, bottom=2mm,
  halign=center
]
  {\Large\bfseries\color{teal!90!black}\faBookmark\ BÀI 58: SO SÁNH HAI PHÂN SỐ}\\[1.5pt]
  {\footnotesize\itshape\color{black!75} 3 trường hợp so sánh cốt lõi: cùng mẫu, khác mẫu, cùng tử và so sánh qua đơn vị trên 1 trang A4}
\end{tcolorbox}

\vspace{-2.5mm}

\noindent
\begin{minipage}[t]{0.49\textwidth}\vspace{0pt}
\begin{tcolorbox}[
  enhanced,
  colback=blue!4!white,
  colframe=blue!75!black,
  arc=2mm, boxrule=0.8pt,
  left=2mm, right=2mm, top=1.8mm, bottom=1.8mm,
  title={\footnotesize\bfseries\color{white}\faSliders\ 1. So sánh cùng mẫu \& Khác mẫu}
]
{\footnotesize
\begin{itemize}\setlength{\itemsep}{1pt}\setlength{\leftskip}{-2.5mm}
  \item \textbf{Cùng mẫu số:} Phân số nào có \textbf{tử số lớn hơn thì lớn hơn}:
  $$\text{Nếu } a > c \implies \dfrac{a}{m} > \dfrac{c}{m}$$
  \item \textbf{Khác mẫu số (2 bước):}
  \begin{itemize}\setlength{\itemsep}{0.5pt}\setlength{\leftskip}{-2.5mm}
    \item Bước 1: Quy đồng mẫu số hai phân số.
    \item Bước 2: So sánh tử số của hai phân số sau khi quy đồng.
  \end{itemize}
\end{itemize}
}
\end{tcolorbox}
\end{minipage}
\hfill
\begin{minipage}[t]{0.49\textwidth}\vspace{0pt}
\begin{tcolorbox}[
  enhanced,
  colback=teal!4!white,
  colframe=teal!75!black,
  arc=2mm, boxrule=0.8pt,
  left=2mm, right=2mm, top=1.8mm, bottom=1.8mm,
  title={\footnotesize\bfseries\color{white}\faRocket\ 2. Kĩ thuật so sánh đặc biệt}
]
{\footnotesize
\begin{itemize}\setlength{\itemsep}{1pt}\setlength{\leftskip}{-2.5mm}
  \item \textbf{Cùng tử số (Quy tắc ngược):} Phân số nào có \textbf{mẫu số bé hơn thì phân số đó LỚN HƠN}:
  $$\text{Nếu } b < d \implies \dfrac{a}{b} > \dfrac{a}{d} \quad (\text{Ví dụ: } \dfrac{3}{4} > \dfrac{3}{5})$$
  \item \textbf{So sánh qua số 1:}\\
  Nếu $\dfrac{a}{b} < 1$ và $\dfrac{c}{d} > 1$ thì kết luận ngay: $\dfrac{a}{b} < \dfrac{c}{d}$.
\end{itemize}
}
\end{tcolorbox}
\end{minipage}

\vspace{-2.5mm}
% --- KHỐI 3: SƠ ĐỒ TƯ DUY TIKZ ---
\noindent
\begin{tcolorbox}[
  enhanced,
  colback=orange!4!white,
  colframe=orange!80!black,
  arc=2mm, boxrule=0.8pt,
  left=2mm, right=2mm, top=1.5mm, bottom=1.5mm,
  title={\footnotesize\bfseries\color{white}\faSitemap\ SƠ ĐỒ TƯ DUY: 3 NGÃ RẼ SO SÁNH PHÂN SỐ}
]
\centering
\begin{tikzpicture}[
  box/.style={draw, rounded corners=2pt, align=center, font=\scriptsize\bfseries, inner sep=3pt, minimum height=1.05cm},
  arr/.style={-{Stealth[length=2.5mm]}, thick}
]
  \node[box, fill=purple!12, draw=purple!80!black, text width=3.0cm] (root) at (0, 0) {SO SÁNH\\HAI PHÂN SỐ};

  \node[box, fill=blue!10, draw=blue!80!black, text width=3.8cm] (b1) at (5.2, 1.15) {CÙNG MẪU SỐ\\Tử lớn hơn $\implies$ Phân số LỚN hơn};

  \node[box, fill=teal!10, draw=teal!80!black, text width=3.8cm] (b2) at (5.2, 0.0) {KHÁC MẪU SỐ\\Quy đồng mẫu $\implies$ So sánh tử};

  \node[box, fill=orange!12, draw=orange!80!black, text width=3.8cm] (b3) at (5.2, -1.15) {CÙNG TỬ SỐ (Ngược lại)\\Mẫu BÉ hơn $\implies$ Phân số LỚN hơn};

  \draw[arr, blue!80!black] (root.east) -- ++(0.4, 0) |- (b1.west);
  \draw[arr, teal!80!black] (root.east) -- (b2.west);
  \draw[arr, orange!80!black] (root.east) -- ++(0.4, 0) |- (b3.west);
\end{tikzpicture}
\end{tcolorbox}

\vspace{-2.5mm}
% --- KHỐI 4: VÍ DỤ & MẸO TRÁNH BẪY ---
\noindent
\begin{minipage}[t]{0.49\textwidth}\vspace{0pt}
\begin{tcolorbox}[
  enhanced,
  colback=green!4!white,
  colframe=green!60!black,
  arc=2mm, boxrule=0.8pt,
  left=2mm, right=2mm, top=1.8mm, bottom=1.8mm,
  title={\footnotesize\bfseries\color{white}\faCheckCircle\ Ví dụ thực hành chuẩn mực}
]
{\footnotesize
\textbf{1. Cùng mẫu:} $\dfrac{3}{7} < \dfrac{5}{7}$ (vì $3 < 5$).\\
\textbf{2. Khác mẫu:} So sánh $\dfrac{3}{4}$ và $\dfrac{5}{6}$ (MSC = 12):\\
$\dfrac{3}{4} = \dfrac{9}{12}; \quad \dfrac{5}{6} = \dfrac{10}{12}$. Vì $\dfrac{9}{12} < \dfrac{10}{12} \implies \mathbf{\dfrac{3}{4} < \dfrac{5}{6}}$.\\
\textbf{3. Cùng tử:} $\dfrac{4}{5} > \dfrac{4}{7}$ (vì $5 < 7$).\\
\textbf{4. Qua số 1:} $\dfrac{2}{3} < 1 < \dfrac{5}{4} \implies \mathbf{\dfrac{2}{3} < \dfrac{5}{4}}$.
}
\end{tcolorbox}
\end{minipage}
\hfill
\begin{minipage}[t]{0.49\textwidth}\vspace{0pt}
\begin{tcolorbox}[
  enhanced,
  colback=red!4!white,
  colframe=red!70!black,
  arc=2mm, boxrule=0.8pt,
  left=2mm, right=2mm, top=1.8mm, bottom=1.8mm,
  title={\footnotesize\bfseries\color{white}\faExclamationTriangle\ Mẹo tránh bẫy trắc nghiệm}
]
{\footnotesize
\begin{itemize}\setlength{\itemsep}{1.5pt}\setlength{\leftskip}{-2.5mm}
  \item \textbf{Bẫy so sánh hai phân số cùng tử:}\\
  Rất nhiều học sinh nhầm: "mẫu lớn hơn thì lớn hơn". Nhớ: Cùng tử thì mẫu BÉ hơn sẽ LỚN hơn (chia bánh cho ít người thì mỗi phần to hơn)!
  \item \textbf{Bẫy chưa quy đồng đã vội so sánh tử.}\\
  Hai phân số khác mẫu bắt buộc phải đưa về cùng mẫu mới được so sánh.
\end{itemize}
}
\end{tcolorbox}
\end{minipage}

\end{document}
